\section{Basis Encoding of Subsets}

Based on the concept of one-hot encodings of states we in this chapter develop the construction of encodings to sets, relations and functions.
We start with the definition of subset encodings, which represent set memberships in their boolean coordinates.

\begin{definition}[Basis encoding of subsets]
    \label{def:subsetEncoding}
    We say that an arbitrary finite set $\arbset$ is enumerated by an enumeration variable $\indvariableof{\arbset}$ taking values in $[\inddimof{\arbset}]$, when $\inddimof{\arbset}=\absof{\arbset}$ and there is a bijective index interpretation function
    \begin{align*}
        \indexinterpretation \defcols [\inddimof{\arbset}] \rightarrow \arbset \, .
    \end{align*}
    Given an set $\arbset$ enumerated by the variable $\indvariableof{\arbset}$, any subset $\arbsubset\subset\arbset$ is encoded by the tensor $\onehotmapto{\arbsubset}[\indvariable]$ defined for $\indindex\in[\absof{\arbset}]$ as
    \begin{align*}
        \onehotmapofat{\arbsubset}{\indexedindvariable}
        = \begin{cases}
              1 & \ifspace \indexinterpretationat{\indindex} \in \arbsubset \\
              0 & \text{else}
        \end{cases} \, .
    \end{align*}
\end{definition}

% Decomposition
In a one-hot basis decomposition we have
\begin{align*}
    \onehotmapofat{\arbsubset}{\indvariable}
    \coloneqq \sum_{\indindex\in[\inddimof{\arbset}]\,:\,\indexinterpretationat{\indindex}\in\arbsubset}\onehotmapofat{\indindex}{\indvariable} \, .
\end{align*}


The inclusion of subsets is represented by the partial ordering of tensors.
Let us first define this property for arbitrary tensors.

% Here for general tensors, not just propositional formulas!
\begin{definition}[Partial ordering of tensors]
    \label{def:partialOrder}
    We say that two tensors $\exformulaat{\shortcatvariables}$ and $\secexformulaat{\shortcatvariables}$ attached with the same variables are partially ordered, denoted by
    \begin{align*}
    {\exformula}
        \prec{\secexformula} \, ,
    \end{align*}
    if for all $\shortcatindices\in\facstates$
    \begin{align*}
        \exformulaat{\indexedshortcatvariables} \leq \secexformulaat{\shortcatindices}  \, .
    \end{align*}
\end{definition}

For boolean tensors, the partially ordering is equal to a subset relation of the coordinates with value $1$, as we show next.

\begin{theorem}
    \label{the:subsetRelationSubsetEncoding}
    Let $\arbset$ be an arbitrary set enumerated by the variable $\indvariable$ and index interpretation function $\indexinterpretation$.
    For two subsets $\arbsetof{0},\arbsetof{1}$ of $\arbset$ we have
    \begin{align*}
        \arbsetof{0} \subset \arbsetof{1}
    \end{align*}
    if and only if
    \begin{align*}
        \onehotmapofat{\arbsetof{0}}{\indvariable} \prec \onehotmapofat{\arbsetof{1}}{\indvariable} \, .
    \end{align*}
\end{theorem}
\begin{proof}
    We have $\arbsetof{0} \subset \arbsetof{1}$ if and only if
    \begin{align*}
        \forall{\indindex\in[\cardof{\arbset}]} \big(\indexinterpretationat{\indindex}\in\arbsetof{0}\big) \Rightarrow \big(\indexinterpretationat{\indindex}\in\arbsetof{1}\big) \, ,
    \end{align*}
    which is equal to
    \begin{align*}
        \forall{\indindex\in[\cardof{\arbset}]} \big(\onehotmapofat{\arbsetof{0}}{\indexedindvariable}=1\big) \Rightarrow \big(\onehotmapofat{\arbsetof{0}}{\indexedindvariable}=1\big)  \, .
    \end{align*}
    Since subset encodings are boolean tensors, this is equivalent to
    \begin{align*}
        \onehotmapofat{\arbsetof{0}}{\indvariable} \prec \onehotmapofat{\arbsetof{1}}{\indvariable} \, . & \qedhere
    \end{align*}
\end{proof}

\subsection{Binary Relations}

% Explanation
%Encoding of subsets as vectors: Each coordinate associated with a possible element, $\{0,1\}$ encoding whether in subset.
%The encodings is thus a boolean tensor.
%Any subset encoding is a boolean tensor.

% Relation
Since relations are subsets of cartesian products between two sets, their encoding is a straightforward generalization of \defref{def:subsetEncoding}.

\begin{definition}[Basis encoding of binary relations]
    A relation between two finite sets $\inset$ and $\outset$ is a subset of their cartesian product
    \begin{align*}
        \exrelation \subset \inset \times \outset \, .
    \end{align*}
    Given an enumeration of $\inset$ and $\outset$ by the categorical variables $\indvariableof{\insymbol}$ and $\indvariableof{\outsymbol}$ and interpretation maps $\indexinterpretationof{\insymbol}$, $\indexinterpretationof{\outsymbol}$, we define the basis encoding of this subset as the tensor $\onehotmapto{\exrelation}[\indvariableof{\insymbol},\indvariableof{\outsymbol}]$ with the coordinates
    \begin{align*}
        \onehotmapofat{\exrelation}{\indexedindvariableof{\insymbol},\indexedindvariableof{\outsymbol}}
        = \begin{cases}
              1 & \ifspace (\indexinterpretationofat{\insymbol}{\indindexof{\insymbol}},\indexinterpretationofat{\outsymbol}{\indindexof{\outsymbol}}) \in \exrelation \\
              0 & \text{else}
        \end{cases} \, .
    \end{align*}
\end{definition}

% Decomposition
The basis encoding has a decomposition into one-hot encodings as
\begin{align*}
    \onehotmapofat{\exrelation}{\indvariableof{\insymbol},\indvariableof{\outsymbol}}
    = \sum_{\indindexof{\insymbol},\indindexof{\outsymbol} \wcols (\indexinterpretationofat{\insymbol}{\indindexof{\insymbol}},\indexinterpretationofat{\outsymbol}{\indindexof{\outsymbol}}) \in \exrelation}
    \onehotmapofat{\indindexof{\insymbol}}{\indvariableof{\insymbol}}  \otimes \onehotmapofat{\indindexof{\outsymbol}}{\indvariableof{\outsymbol}}  \, .
\end{align*}

basis encodings have a matrix structure by the cartesian product, which can be further folded to tensors, when the sets itself are cartesian products.
The basis encoding is a bijection between the relations of two sets and the boolean tensors with their enumeration variables.

%They provide representations of generic relations by boolean tensors, in the sense that each relation between two sets is represented
%\begin{theorem}
%	The basis encoding is a bijection between the set of relations and the set of boolean tensors.
%\end{theorem}
%\begin{proof}
%	% =>
%	By definition, a basis encoding is the encoding of a subset and thus a boolean tensor.
%	% <=
%	Any matrification of a boolean tensor marks by its $1$ coordinates the elements of a relation.
%\end{proof}
%
%% Significance
%We can thus understand any matrification of a boolean tensor as the encoding of a relation and vice versa.



\subsection{Higher-Order Relations}

We can extend this contraction to relations of higher order, and arrive at encoding schemes usable for relational databases.

\begin{definition}[Basis encoding of $\atomorder$-ary relations]
    \label{def:daryRelation}
    Given sets $\arbsetof{\atomenumerator}$ for $\atomenumeratorin$, a $\atomorder$-ary relation is a subset of a their cartesian product, that is
    \begin{align*}
        \exrelation \subset\bigtimes_{\atomenumeratorin} \arbsetof{\atomenumerator} \, .
    \end{align*}
    Given an enumeration of each set $\arbsetof{\atomenumerator}$ by a variable $\indvariableof{\atomenumerator}$ and an interpretation map $\indexinterpretationof{\atomenumerator}$, we define the basis encoding of the relation as the tensor $\onehotmapto{\exrelation}[\indvariableof{[\atomorder]}]$ with coordinates
    \begin{align*}
        \onehotmapofat{\exrelation}{\indexedindvariableof{[\catorder]}}
        = \begin{cases}
              1 & \ifspace (\indexinterpretationofat{0}{\indindexof{0}},\ldots,\indexinterpretationofat{\atomorder-1}{\indindexof{\atomorder-1}}) \in \exrelation \\
              0 & \text{else}
        \end{cases} \, .
    \end{align*}
\end{definition}

%\begin{example}[Propositional Formulas]
Let there be for $\atomenumeratorin$ sets $\arbsetof{\atomenumerator}$ of truth assignments to the $\atomenumerator$-th atom, which are all enumerated by $[2]$.
A propositional formula then corresponds with a $\atomorder$-ary relation and we directly defined them in \defref{def:formulas} by their basis encoding.
%\end{example}

% Minterm interpretation
%If we demand $\catdimof{\atomenumerator}=2$, we can interpret the one-hot encoding of  propositional formulas.

% Knowledge Bases
\begin{theorem}
    The encoding of any $\catorder$-ary relation
    \begin{align*}
        \exrelation = \{ \shortcatindices^{\decindex} \wcols \decindexin \} \subset \bigtimes_{\catenumeratorin} \truthset \,
    \end{align*}
    where the objects in $\truthset$ are enumerated by $\catvariableof{\atomenumerator}$ with the standard index interpretation function (see \secref{sec:booleanEncoding})
    \begin{align*}
        \indexinterpretationat{\truesymbol} = 1 \andspace \indexinterpretationat{\falsesymbol} = 0 \, ,
    \end{align*}
    coincides with the propositional formula
    \begin{align*}
        \formulaat{\shortcatvariables} = \bigvee_{\decindexin} \termof{\catindex_{[\atomorder]}^{\decindex}} \, .
    \end{align*}
\end{theorem}
\begin{proof}
    By definition, the encoding $\onehotmapof{\exrelation}$ is decomposed as
    \begin{align*}
        \onehotmapofat{\exrelation}{\shortcatvariables}
        = \sum_{\decindexin} \onehotmapofat{\shortcatindices^{\decindex}}{\shortcatvariables} \, .
    \end{align*}
    By \theref{the:tensorToMaxMinTerms} this is equal to
    \begin{align*}
        \hypercoreat{\shortcatvariables} = \left( \bigvee_{\hyperonecoordinates}
        \termof{\catzeropositions}{\catonepositions}
        \right)[\shortcatvariables]
        = \bigvee_{\decindexin} \termof{\catindex_{[\atomorder]}^{\decindex}} \, . & \qedhere
    \end{align*}
\end{proof}


\begin{example}[Relational Databases]
    Relational Databases can be encoded as tensors using the relation encoding scheme.
    Each column is thereby understood as an enumeration variable, which values form the sets $\arbsetof{\catenumerator}$.
\end{example}

% Sparse Representations
Let us notice, that the dimensionality of the tensor space used for representing a relation is
\begin{align*}
    \prod_{\catenumeratorin} \cardof{\arbsetof{\catenumerator}}
\end{align*}
and therefore growing exponentially with the number of variables.
Relations are however often sparse, in the sense that
\begin{align*}
    \cardof{\exrelation} << \prod_{\catenumeratorin} \cardof{\arbsetof{\catenumerator}} \, .
\end{align*}
It is therefore often beneficially to choose sparse encoding schemes, for example by restricted CP formats (see \charef{cha:sparseRepresentation}) to represent $\onehotmapof{\exrelation}$.

